\section*{Построение арбитражного портфеля в однопериодической модели}

\subsection*{Условие}

Даны:
\begin{itemize}
    \item Цена безрискового актива: $B_0 = 10$
    \item Цена рисковых активов: $S_1^0 = 55$, $S_2^0 = 22$
    \item Безрисковая ставка: $r = 10\%$
    \item Ограничение на заём: $1000$
\end{itemize}

Дивиденды:
\begin{itemize}
    \item Благоприятный исход:
    \[
    d_1 = 20\%, \quad d_2 = 50\%
    \]
    \item Неблагоприятный исход:
    \begin{itemize}
        \item первая акция: половина дохода
        \item вторая акция: треть дохода
    \end{itemize}
\end{itemize}

\subsection*{Динамика активов}

Безрисковый актив:
\[
B_1 = B_0(1+r) = 10 \cdot 1.1 = 11
\]

\paragraph{Благоприятный исход:}
\[
S_1^U = 55(1+0.2) = 66
\]
\[
S_2^U = 22(1+0.5) = 33
\]

\paragraph{Неблагоприятный исход:}
\[
S_1^D = 55\left(1+\frac{0.2}{2}\right) = 55 \cdot 1.1 = 60.5
\]
\[
S_2^D = 22\left(1+\frac{0.5}{3}\right) = 22 \cdot \frac{7}{6} \approx 25.67
\]

\subsection*{Портфель}

Рассмотрим портфель:
\[
X = \alpha B + \beta_1 S_1 + \beta_2 S_2
\]

Начальная стоимость:
\[
X_0 = 10\alpha + 55\beta_1 + 22\beta_2
\]

Стоимость в момент времени $t=1$:
\[
X_U = 11\alpha + 66\beta_1 + 33\beta_2
\]
\[
X_D = 11\alpha + 60.5\beta_1 + 25.67\beta_2
\]

\subsection*{Условия арбитража}

Арбитраж существует, если:
\[
X_0 \le 0, \quad X_U \ge 0, \quad X_D \ge 0
\]
и хотя бы одно неравенство строгое.

\subsection*{Построение портфеля}

Рассмотрим выбор:
\[
\beta_1 = -1, \quad \beta_2 = 1
\]

Найдём $\alpha$ из условия $X_0 = 0$:
\[
10\alpha - 55 + 22 = 0
\]
\[
10\alpha = 33 \Rightarrow \alpha = 3.3
\]

\subsection*{Проверка}

Начальная стоимость:
\[
X_0 = 10 \cdot 3.3 - 55 + 22 = 0
\]

\paragraph{Благоприятный исход:}
\[
X_U = 11 \cdot 3.3 - 66 + 33 = 36.3 - 33 = 3.3 > 0
\]

\paragraph{Неблагоприятный исход:}
\[
X_D = 11 \cdot 3.3 - 60.5 + 25.67 = 36.3 - 34.83 = 1.47 > 0
\]

\subsection*{Вывод}

Получен арбитражный портфель:
\[
\alpha = 3.3, \quad \beta_1 = -1, \quad \beta_2 = 1
\]

Он обладает следующими свойствами:
\begin{itemize}
    \item не требует начальных вложений;
    \item приносит положительную прибыль в любом состоянии рынка.
\end{itemize}

Следовательно, на рынке существует арбитражная возможность.
